\section*{Appendix B: Event Scheduling (Workplan)}

This appendix presents the project workplan outlining the planned activities for the development of the Traxis system. The workplan is organized chronologically and aligns with the Agile-inspired methodology described in Chapter 4. It defines the scheduled development events, their durations, and intended outputs, providing a structured roadmap for the successful execution of the project.

\subsection*{Month 1: Foundation Phase (January 2026)}

\textbf{Activity 1: Backend and Database Setup} \\
\textit{Scheduled Period: 1 January 2026 – 14 January 2026} \\
During this period, the project team will establish the foundational backend infrastructure for the system. This will include configuring the Django backend framework, setting up the PostgreSQL database with PostGIS extensions, defining core data models, and creating initial RESTful API endpoints. The expected outcome will be a functional backend environment capable of supporting authentication, data persistence, and future feature development.

\textbf{Activity 2: Authentication and Location Services Implementation} \\
\textit{Scheduled Period: 15 January 2026 – 31 January 2026} \\
This activity will focus on implementing secure authentication mechanisms and geolocation services. Firebase Authentication will be integrated to support phone-number-based user login, while Google Maps APIs will be configured for real-time location tracking and routing. Backend logic for managing user sessions and location updates will be implemented to support subsequent ride-matching functionality.

\subsection*{Month 2: Core Interaction Phase (February 2026)}

\textbf{Activity 3: User Interface and Authentication Workflow Development} \\
\textit{Scheduled Period: 1 February 2026 – 14 February 2026} \\
During this phase, the client-side applications will be developed using Flutter. User interface components, navigation flows, and authentication workflows for both passengers and drivers will be implemented from a shared codebase. The primary objective will be to deliver a functional, responsive, and user-friendly mobile interface capable of interacting securely with backend services.

\textbf{Activity 4: Ride Management Logic and AI Integration} \\
\textit{Scheduled Period: 15 February 2026 – 28 February 2026} \\
This activity will involve implementing the core ride management features of the system. Ride request handling, driver availability broadcasting, and passenger–driver matching logic will be developed within the backend. AI services will be integrated through MCP to support predictive matching and demand-aware allocation. Real-time notifications will be enabled using Firebase Cloud Messaging to ensure timely communication between system components.

\subsection*{Month 3: System Enhancement and Finalization Phase (March 2026)}

\textbf{Activity 5: Analytics and Administrative Tools Development} \\
\textit{Scheduled Period: 1 March 2026 – 14 March 2026} \\
This phase will focus on the development of administrative and analytical capabilities. The Admin Portal will be implemented using Vue.js to provide tools for user management, trip monitoring, and system analytics. Supporting backend endpoints for reporting, logging, and administrative operations will be completed to enable effective system oversight.

\textbf{Activity 6: Performance Optimization, Accessibility, and System Validation} \\
\textit{Scheduled Period: 15 March 2026 – 31 March 2026} \\
The final activity will concentrate on system refinement and deployment readiness. Performance optimization will be conducted across the backend and mobile applications, including database tuning and background task optimization. Accessibility enhancements will be implemented to improve usability for diverse user groups. Comprehensive testing, bug fixing, and final system validation will be carried out to ensure compliance with project requirements.

\subsection*{Workplan Summary}

The proposed workplan will ensure a phased and systematic approach to system development, progressing from foundational setup to core functionality and final optimization. By adhering to clearly defined activities and timelines, the project will maintain effective coordination, reduce development risks, and support the timely delivery of the Traxis system in accordance with the proposed methodology.